%! Author = Maximilian Fülle
%! Date = 29.05.2026


\chapter{Valued Fields and Models}\label{ch:valued-fields-and-models}

\section{Models of Curves}\label{sec:models-of-curves}

\subsection{Basic Definitions}\label{subsec:basic-definitions}
Let $K$ be a field endowed with an additive valuation of rank $1$, $v_K: K \to \hat{\RR} \coloneqq \RR\cup \lbrace\infty\rbrace$.
Let $\Gamma_K \coloneqq v_K(K^*) \subset \RR$ be the value group of $v_K$, $\OK \subset K$ its valuation ring with
maximal ideal $\mK$ and $k \coloneqq \OK /\mK$ its residue field.

For most of the the first chapter, we make the following additional assumptions:
\begin{assumption} \label{ass:v_K}
    \begin{enumerate}[(a)]
        \item $v_K$ is discrete, i.e.\ $\Gamma_K$ is a discrete subgroup of $\RR$,
        \item $K$ is complete with respect to $v_K$,
        \item $k$ is perfect.
    \end{enumerate}
\end{assumption}

Throughout the whole monograph, $X$ will denote a \textit{nice} curve over $K$,
that is a smooth, projective and absolutely irreducible curve over $K$.

\begin{definition} \label{def:model}
    A \textit{model} $\X$ of $X$ is a proper, flat and normal $\OK $-scheme
    with an identification of its generic fibre $\X_g$ with $X$.
\end{definition}

\begin{example}
    Let $K = \QQ_p$, $\OK  = \ZZ_p$ and $X = \PP^1(K)$, then $\X = \PP^1(\OK )$ is a model of $X$ with the
    evident isomorphism from $\X_g = \X \times_{\OK} K$ to $\PP^1(K)$.
    The special fibre\footnote{That is, the fibre of the closed point of $\Spec \OK $.}
    $\X_s$ has one irreducible component which is isomorphic to $\PP^1(k)$.
\end{example}

\begin{remark} \label{rem:flatIffIntegral}
    The condition of a model being flat is equivalent to it being integral:
    If $\X \to \Spec \OK $ is an integral scheme, then for every $x\in \X$, the stalk $\OO_{\X,x}$ is an integral ring extension of $\OK$.
    Thus $\OO_{\X,x}$ is a torsion-free module over the discrete valuation ring $\OK$, hence flat by~\cite[Corollary 1.2.14]{Liu}.
    Conversely, the integrality of a model $\X$ of $X$ follows from its flatness and the integrality of $X$ by~\cite[Proposition 4.3.8]{Liu}.
\end{remark}

\begin{remark}
    A point $x\in \X$ of codimension $1$ which is regular is normal.
    This is conversely true by Serre's criterion (see~\cite[Lemma 8.2.21]{Liu}).
\end{remark}

\begin{example}\label{ex:model}
    Consider the curve $X$ defined by $y^2 = x^3 + px$ over $\QQ_p$ with $p \neq 2$.
    We show that
    \begin{align}
        \X \coloneqq \Proj \ZZ_p[X,Y,Z]/(Y^2 Z - X^3 - pXZ^2)
    \end{align}
    is a model of $X$.
    It is clearly a proper scheme over $\OK  = \ZZ_p$.
    For the flatness it suffices to show that
    $R\coloneqq \ZZ_p[X,Y,Z]/(Y^2 Z - X^3 - pXZ^2)$ is a torsion-free $\ZZ_p$-module by Remark~\ref{rem:flatIffIntegral}.
    But this is true since $Y^2 Z - X^3 - pXZ^2$ is irreducible in $\ZZ_p[X,Y,Z]$. \\
    For normality, consider the affine patch $U: \, Z \neq 0$, that is
    \begin{align}
        U \coloneqq \Spec{\ZZ_p[x,y]/(y^2 - x^3 + px)}
    \end{align}
    where $x = X/Z$ and $y = Y/Z$.
    The special fibre $\X_s$ is defined by the equation $\bar{X}_p: \, \bar{y}^2 - \bar{x}^3 = 0$, which is irreducible in $\FF[x,y]$.
    The only singularity is when $\bar{x}=\bar{y}=0$.
    Since $\X$ is regular at all other points, condition $(\mathrm{S}_2)$ 
    of Serre's criterion for normality holds. Since $\X_s$ has no embedded points,
    $(\mathrm{R}_1)$ is also true, which means that $\X$ is normal (cf. Lemma~\ref{lem:conditions-for-normality}).
    This shows that $\X$ is indeed a model of $X$.
    %Let $Z$ be the irreducible component of $\X_s$ containing the generic point $\eta$.
    %The only singularity is when $\eta:\, y = x = 0$ and so one has the maximal ideal $\m_p = (x,y,p)$ of $\OO_{Z, \eta}$.
    %All monomials have degrees greater than $2$, so $\m_p$ can't be generated by less than three elements modulo $\m_p^2$.
    %This implies
    %\begin{align}
    %    \dim{\FF_p}{\m_p/\m_p^2} = 3 > 2 = \codim{\eta},
    %\end{align}
    %i.e. $\eta$ is a non-regular point of $U$.
\end{example}


\subsection{From Models to Valuations} \label{subsec:from-valuations-to-models}
In this section, we assume the valuation $v_K$ to be \textit{normalised}, i.e. $v_K(K^*) = \ZZ$.

\begin{definition}
    Let $F_X$ be a function field of $X$ over $K$.
    A valuation $v: F_X \to \hat{\RR}$ is said to be \textit{of type \textrm{II}}, if
    \begin{enumerate}
        \item $v|_K = v_K$,
        \item $v$ is discrete, and
        \item the extension of residue fields $k(v)|k$ has transcendence degree $1$.
    \end{enumerate}
    A valuation of this type is also known as \textit{geometric valuation}.
    We write $V_\textrm{II}(F_X)$ for the set of all valuations of type \textrm{II} of $F_X$.
    The number
    \begin{align}
        m(v) \coloneqq [\Gamma_v : \Gamma_K]
    \end{align}
    is called the \textit{multiplicity of $v$}, where $\Gamma_v \coloneqq v(F_X^*)$ is the value group.
    We say that $v$ is \textit{weakly unramified}, if $m(v) = 1$.
\end{definition}

\begin{example}
    The \textit{Gauß valuation} on the polynomial ring $K[x]$ is defined as
    \begin{align}
        v_{\scriptscriptstyle G}: K[x] &\to \hat{\QQ}, \; v_{\scriptscriptstyle G}\left( \sum_i a_i x^{i} \right) = \min_i(v_K(a_i)).
    \end{align}
    Since $k(v_{\scriptscriptstyle G}) = k(\bar{x})$ (where $\bar{x}$ is the image of $x$ in $k(v_g)$), the Gauß valuation
    extends naturally and uniquely to $K(x)$, meaning that we have $v_{\scriptscriptstyle G} \in V_\textrm{II}(K(x))$.
    We will later see that $\lbrace v_{\scriptscriptstyle G} \rbrace$ defines the model $\X = \PP^1(\ZZ_p)$ of $\PP^1(\QQ_p)$.
\end{example}

Let $\X$ be a model of $X$.
We wish to assign a valuation in $V_\textrm{II}(F_X)$ to each generic point of the special fibre $\X_s$.
Suppose that $Z$ is an irreducible component of $\X_s$ with generic point $\xi$.
Since $\xi$ has codimension $1$ in $\X$, the stalk $\OO_{\X,\xi}$ is a discrete valuation ring\@.
Thus, we can define a valuation on $F_X$,
\begin{align} \label{eq:ValuationOfComponent}
    v_Z: \mathrm{Frac}(\OO_{\X, \xi}) = F_X \to \hat{\QQ}.
\end{align}
We set
\begin{align}
    V(\X) \coloneqq \lbrace v_Z: \; Z \subset \X_s \text{ is an irreducible component}\rbrace.
\end{align}

\begin{proposition}\label{prop:WellDefOfInducedValuation}
    The valuation defined in~\eqref{eq:ValuationOfComponent} is of type \textrm{II}.
\end{proposition}

\begin{proof}
    Since $\X$ is a flat $\OK $-scheme $\OO_{\X, \xi}$ dominates $\OK$, and we may scale $v_Z$ such
    that the restriction of $v_Z$ to $K$ is $v_K$.
    Furthermore, we have a closed immersion $Z \hookrightarrow \X$
    by definition, inducing a surjection $\OO_{\X,\xi} \to \OO_{Z,\xi}$.
    Thus $\trdeg(k(v_Z)) = \codim(\OO_{Z,\xi}) = \codim(v_Z) = 1$, showing $v_Z \in V_\textrm{II}(F_X)$.
\end{proof}

\begin{definition}
    The \textit{multiplicity} of the irreducible component $Z$ of $\X_s$ containing the generic point $\xi$ is defined as
    \begin{align}
        m(Z) \coloneqq \length(\OO_{\X,\xi}).
    \end{align}
\end{definition}

\begin{proposition}\label{prop:multiplicities-coincide}
    It holds $m(v_Z) = m(Z)$.
\end{proposition}

\begin{proof}
    As the model $\X$ is normal, the stalk $\OO_{\X, \xi}$ is a discrete valuation ring\@.
    Let $\pi_\xi$ denote an uniformizing parameter, then
    \begin{align}
        \OO_{\X_s,\xi} = \OO_{\X,\xi}/(\pi_\xi).
    \end{align}
    Thus, the length of the stalk at $\xi$ coincides with the length of the chain of ideals generated
    by the uniformizing parameter of the quotient ring:
    Since $(\pi_\xi) = (\pi_{Z,\xi}^m)$, where $m = (\Gamma_{v_Z}:\Gamma_K)$, it holds $m(v_Z) = \length(\OO_{\X, \xi})$ as desired.
\end{proof}

\begin{definition}
    A \textit{modification} of a model of $X$ is a morphism $\X' \to \X$ of $\OK $-schemes which induces the
    identity on the generic fibre.
\end{definition}

\begin{remark}\label{rem:ModificationsAreSurjective}
    Since a model is proper over $\OK $, a modification is a proper morphism.
    Thus since it is dominant, it is also surjective.
\end{remark}

The main goal of the first two sections of this chapter is to establish the following theorem.

\begin{theorem}\label{thm:bijectionModelsAndComponents}
    The map
    \begin{align} \label{eq:mapFromModelToValuation}
        \X \mapsto V(\X).
    \end{align}
    defines a monotone bijection between
    \begin{itemize}
        \item the set of isomorphism classes of models of $X$, and
        \item the set of finite, nonempty subsets of $V_\textrm{II}(F_X)$.
    \end{itemize}
\end{theorem}

In the remainder of this section, we prove injectivity of this map.

\begin{lemma}[{\cite[Lemma 3.6]{Rueth}}]\label{lem:preimageOfGenericPoint}
    Let $\X$ be a model of $X$ and let $\X'$ be a model of another irreducible smooth projective curve $X'$.
    Let $\varphi: \X' \to \X$ be a morphism of $\OK $-schemes which induces a cover $X' \to X$ on the generic fibre.
    Let $\eta$ be a generic point of $\X_s$, then the preimage of $\eta$ under $\varphi$ is non-empty and consists of generic points of $\X_s'$.
    If $\varphi$ is birational, then the preimage consists of precisely one generic point.
\end{lemma}

\begin{proof}
    Let $Z \subset \X_s$ be the irreducible component that contains $\eta$.
    Since $\varphi$ is surjective (cf.\ Remark~\ref{rem:ModificationsAreSurjective}),
    there is a point $\eta'$ on an irreducible component $Z' \subset \X_s'$ which is mapped to $\eta$.
    Since $\varphi$ is dominant, the induced ring homomorphism $\OO_{\X,\eta} \to \OO_{\X',\eta'}$ is injective.
    This is an extension of discrete valuation rings and thus $\eta' \in Z'$ must be generic.
    In particular, if $\varphi$ is birational, then both rings are integrally closed in the same field of fractions which makes them equal.
    Therefore, $\eta'$ must be unique since $\X'$ is separated.
\end{proof}

\begin{lemma}[{\cite[Lemma 3.7]{Rueth}}]
    Let $\varphi: \X' \to \X$ be a modification of models of $X$.
    Then $V(\X) \subset V(\X')$.
\end{lemma}

\begin{proof}
    This is a consequence of Proposition~\ref{prop:WellDefOfInducedValuation} and Lemma~\ref{lem:preimageOfGenericPoint}.
\end{proof}

We have a partial order of models of $X$ given by $\X \leq \X'$ if there exists a modification $\X' \to \X$.
By the previous lemma, this gives a partial order on the corresponding finite subsets of $V_\textrm{II}(F_X)$ via inclusion maps.

\begin{lemma}[{\cite[Lemma 3.8]{Rueth}}]
    Let $\X$ and $\X'$ be models of $X$ with $V(\X) \subset V(\X')$.
    Then there is a modification $\X' \to \X$.
\end{lemma}

\begin{proof}[Sketch of the proof]
    Let
    \begin{align*}
        U= \lbrace Z_\eta \subset \X_s: v_{Z_\eta} \in V(\X)\rbrace
    \end{align*}
    and
    \begin{align*}
        U'= \lbrace Z_{\eta'} \subset \X_s: v_{Z_{\eta'}} \in V(\X')\rbrace
    \end{align*}
    be covers of $\X_s$ (resp. $\X_s'$) of irreducible components, where the irreducible component
    $Z_\eta$ has generic point $\eta$.
    The inclusion $V(\X) \subset V(\X')$ induces an injective map $\iota$ from the generic points of the irreducible components of $\X_s$ to those of $\X'_s$, which induces
    ring homomorphisms $\OO_{\X_s', \iota(\eta)} \to \OO_{\X_s, \eta}$.
    Thus, we get a birational map $\X' \dbkarow \X$ which is a morphism by~\cite[Theorem 8.3.20]{Liu}.
\end{proof}

\begin{proposition}[{\cite[Proposition 3.9]{Rueth},\cite[Corollary 8.3.23]{Liu}}]
    Two models $\X$ and $\X'$ of $X$ are isomorphic as models of $X$ if and only if $V(\X) = V(\X')$.
    In particular, the map defined in~\eqref{eq:mapFromModelToValuation} is injective.
\end{proposition}

\begin{proof}
    Suppose that $V(\X) = V(\X')$.
    Then there exist morphisms $\X\to \X'$ and $\X' \to \X$.
    The composition of these morphisms belongs to the equivalence class of the identity morphism, when considered as a rational map.
\end{proof}

\begin{example}
    We continue with Example~\ref{ex:model} of the model
    \begin{align*}
        \mathcal{X} = \Proj \ZZ_p[X,Y,Z]/(Y^2 Z - X^3 - pXZ^2)
    \end{align*}
    of the elliptic curve $X:\; y^2 = x^3 +px$ defined over $\QQ_p$.
    We have already checked that it is a model.
    Let $Z$ be the irreducible component in $\X_s$ containing the generic point $\xi$ and denote the corresponding valuation by $v_Z$.
    $Z$ is defined by $\Bar{X}_p:\; y^2 = x^3$, thus $Z = \X_s$ and $V(\X) = \lbrace v_Z\rbrace$. \\
    We want to describe $v_Z$ as
    \begin{align*}
        v_Z|_{F_X} = v_{\scriptscriptstyle G}
    \end{align*}
    with respect to the $x$-coordinate.
    Using Lemma~\ref{lem:conditions-for-normality}, we see that $v_Z(p) = 1$
    and $v_Z(x) = 0$, since the horizontal divisor on $x = 0$ intersects $\Bar{X}_p$ at a unique point
\end{example}


\section{From Valuations to Models}\label{sec:from-valuations-to-models}

For the proof of Theorem~\ref{thm:bijectionModelsAndComponents}, it remains to show the surjectivity of the map (\eqref{eq:mapFromModelToValuation}),
i.e., for a given finite, non-empty set of valuations
$V \subset V_\textrm{II}(F_X)$, there is a model $\X$ of $X$ such that $V = V(\X)$.
We will prove this in two steps, which we state in the following proposition.

\begin{proposition}[{{\cite[3.14, 3.16]{Rueth}}}, {{\cite[8.3.36]{Liu}}}] \label{prop:constructing_models}
    \begin{enumerate}[(i)]
        %\item[]
        \item Given model $\X$ and a valuation $v \in V_\textrm{II}(F_X)$, there exists a modification
        $\varphi \colon \X' \to \X$ such that $v \in V(\X')$. \label{prop:constructing_models_i}
        \item Given a model $\X$ and a nonempty subset $V \subset V(\X)$,
        there exists a modification $\varphi \colon \X \to \X'$ such that $V = V(\X')$. \label{prop:constructing_models_ii}
    \end{enumerate}
\end{proposition}

After proving Proposition~\ref{prop:constructing_models}, the proof is complete:
Starting with a model $\X$ of $X$, we can apply part (\ref{prop:constructing_models_i}) for every $v \in V \subset V_\textrm{II}(F_X)$
to obtain a new model that induces all valuations in $V$, but possibly more.
We then apply part (\ref{prop:constructing_models_ii}) to construct a model that induces exactly the valuations in $V$.


\subsection{Proof of the First Part of Theorem~\ref{thm:bijectionModelsAndComponents}}\label{subsec:proof-of-the-first-part}

We fix a model $\X$ of $X$ and a type II valuation $v \in V_\textrm{II}(F_X)$.

\begin{definition}
    A point $x \in \X$ is called a \textit{centre} of $v$ if there exists a morphism $f \colon \Spec \OO_v  \to \X$ of $\OK $-schemes
    which maps the closed point to $x$ and makes the following diagram commute:
    \[
        \begin{tikzcd}
            \Spec \OO_v \arrow[rr, "f"] &                                 & \X \\
            & \Spec F_X \arrow[ru] \arrow[lu] &
        \end{tikzcd}
    \]
    Here, the left map is induced by the inclusion $\OO_v \hookrightarrow \Frac \OO_v = F_X$,
    and the right map comes from the identification of $X$ with the special fibre $\X_s$.
    Specifically, it is the canonical map $\Spec F_X \cong \Spec \OO_{\X, \xi} \to \X$, where $\xi$ is the generic point of $\X$.
\end{definition}

\begin{lemma}
    There exists a unique centre $x \in \X$ of $v$.
\end{lemma}

\begin{proof}
    By the valuative criterion for properness, such a centre exists, and the morphism $f$ is unique
    (see also~\cite[\href{https://stacks.math.columbia.edu/tag/01J5}{Section 01J5, Lemma 01J6}]{Stacks}):
    \[
        % https://tikzcd.yichuanshen.de/#N4Igdg9gJgpgziAXAbVABwnAlgFyxMJZABgBpiBdUkANwEMAbAVxiRAB12BlNGAYwAEAMQD6ADRABfUuky58hFAEZSSqrUYs2nHvwGcA8gZEBpKTJAZseAkTJrq9Zq0QduvQYeM1zs6wqIVSkdNFzcJSXUYKABzeCJQADMAJwgAWyQAJmocCCQlaSTUjMQAZhy8xAKLFPSkMhBcrMKQWpKGprKWtqyKpHKQBiwwMKgIHBxokBDnNkSpCkkgA
        \begin{tikzcd}[column sep=large, row sep=large]
            \Spec F_X \arrow[d] \arrow[r]                 & \X \arrow[d] \\
            \Spec \OO_v \arrow[r] \arrow[ru, "f", dotted] & \Spec \OK  \cdpunkt
        \end{tikzcd}
    \]
\end{proof}

\begin{lemma}[{{\cite[3.11]{Rueth}}}]\label{lem:codim1}
    We have $v \in V(\X)$ if and only if the centre $x$ of $v$ is of codimension one.
    In that case $x$ is the generic point on the special fibre which induces $v$.
\end{lemma}

\begin{proof}
    Assume $v \in V(\X)$.
    Then $v$ is induced by a generic point $\eta \in \X_s$ on the special fibre.
    The canonical morphism $\Spec \OO_{\X, \eta} \to \X$ shows that $v$ has centre $\eta$,  which is indeed a point of codimension one.
    Conversely, suppose $v \in V_\textrm{II}(F_X)$ has center $x \in \X$ of codimension one.
    The morphism $\Spec \OO_v \to \X$ uniquely factors through $\Spec \OO_v \to \Spec \OO_{\X, x}$, so $\OO_v$ dominates $\OO_{\X, x}$.
    Since these are both valuation rings with common fraction field $K$, it follows that they are in fact equal.
    Thus \(v \in V(\X)\).
\end{proof}

\begin{lemma}[{{\cite[3.12, 3.13]{Rueth}}}] \label{lem:condition_induced}
    Let $x \in \X$ be the centre of $v$, and let $f \in \OO_v^* \subset F_X$ be such that the reduction $\overline{f} \in k(v)$ is transcendental over $k$.
    Then:
    \begin{enumerate}[(a)]
    \item If $f \in \OO_{\X,x}$, then $v \in V(\X)$ \label{lem:condition_induced_i}
    \item If the rational map $\X \dashrightarrow \PP^1(\OK )$ induced by $f$ is defined at the center $x$ of $v$,
    then $v \in V(\X)$. \label{lem:condition_induced_ii}
    \end{enumerate}
\end{lemma}

\begin{proof}
    \begin{enumerate}[(a)]
    \item By Lemma~\ref{lem:codim1}, it suffices to show that $x$ is a point of codimension one.
    Choose an open affine set $\Spec A = U \subset \X$ such that $x \in U$ and $f \in A$.
    Denote by $\p \subset A$ the prime ideal corresponding to $x$.
    Since $v$ is non-trivial, the center $x$ is not the generic point of $\X$, and so the codimension is not zero.
    If the dimension of $A/\p$ is at least one, then the inequality
    \[
        \dim (A_\p) + \dim (A/\p) \leq \dim (A) = 2
    \]
    implies
    \[
        \codim(\overline{\{x\}}, \X) = \dim (A_\p) \leq 2 - \dim(A/\p) \leq 2 - 1 = 1,
    \]
    and hence equals one.
    Thus, it suffices to show that $A/\p$ is not a field.
    If $A/\p$ were a field, it would be a finite extension
    of $k$, but this is impossible since $\overline{f} \in A/\p \subset k(v)$ is transcendental over $k$.
    \item It suffices to show, using (\ref{lem:condition_induced_i}), that $f \in \OO_{\X, x}$.
    Let $P = V(f^{-1})$ and $Z = V(f)$ be the sets of poles and zeros of $f$, respectively.
    Then the rational map $\X \dashrightarrow \PP^1(\OK )$ is obtained by gluing the morphisms
    \[
        \X \setminus P  \to \Spec \OK [t] \subset \PP^1(\OK ), \quad \textrm{and} \quad
        \X \setminus Z \to \Spec \OK [t^{-1}] \subset \PP^1(\OK ),
    \]
    which are induced by
    \begin{center}
        \begin{minipage}{0.25\textwidth}
            \begin{align*}
                \OK [t] &\to \Gamma(\X \setminus P, \OO_{\X}), \\
                t &\mapsto f
            \end{align*}
        \end{minipage}
        \quad and \quad
        \begin{minipage}{0.25\textwidth}
            \begin{align*}
                \OK [t^{-1}] &\to \Gamma(\X \setminus Z, \OO_{\X}), \\
                t^{-1} &\mapsto f^{-1}
            \end{align*}
        \end{minipage}
    \end{center}
    and is defined on $U \coloneqq (\X \setminus P) \cup (\X \setminus Z) = \X \setminus (P \cap Z)$.
    By assumption, $x \in U$. If $x \in \X \setminus P$, the claim follows immediately, as
    \[
        f \in \Gamma(\X \setminus P, \OO_{\X}) \subset \OO_{\X, x}.
    \]
    If $x \in \X \setminus Z$, then
    \[
        f^{-1} \in \Gamma(\X \setminus Z, \OO_{\X}) \subset \OO_{\X, x}.
    \]
    The image of $f^{-1}$ under the local ring homomorphism $\OO_{\X, x} \to \OO_v$ is a unit by assumption,
    hence $f^{-1}$ is also a unit in $\OO_{\X, x}$.
    Consequently $f \in \OO_{\X, x}$.
    \end{enumerate}
\end{proof}

We can now finish the proof of Proposition \ref{prop:constructing_models} \ref{prop:constructing_models_i}.

\begin{proof}[Proof of Proposition \ref{prop:constructing_models} \ref{prop:constructing_models_i}]
    Let $f \in \OO_v^\times \subset F_X$ be such that the reduction $\overline{f} \in k(v)$ is transcendental over $k$.
    We also write $f \colon \X \dashrightarrow \PP^1(\OK )$ for the rational map induced by $f$, and let $U$ be its domain of definition.
    Define $\Gamma_f \subset \X \times_{\OK } \PP^1(\OK )$ to be the graph of $f$.
    Recall that it is defined to be the closure of the graph $f|_U$ in $\X \times_{\OK } \PP^1(\OK ) \supset U \times_{\OK } \PP^1(\OK )$.
    This is almost a model of $X$, as it satisfies all the properties of a model except that it might not be normal.
    \begin{center}
        % https://tikzcd.yichuanshen.de/#N4Igdg9gJgpgziAXAbVABwnAlgFyxMJZARgBpiBdUkANwEMAbAVxiRAB12BxOgW17oB9AGYgAvqXSZc+QigAM5KrUYs2nABoACTnl7xBwTgHljggNJid7AAo2AesUMmzl8ZJAZseAkQBMStT0zKyIHLYOTkbsphZi7lLeskRkfsrBamGaCZ7SPnIkpPLpqqHhGgDk4sowUADm8ESgwgBOELxIiiA4EEhk3XRYDGwAFhAQANY5re2d1D1IfhLNbR2IXQuIAMzLIDNrW-O9iAEgDFhgZVAQODi1IEGlbKLUDHQARjAMNnnJYS1YOojHDTVZIAAsR06YgoYiAA
        \begin{tikzcd}
            & \X' \arrow[d]                                &               \\
            \X \times_{\OK } \PP^1(\OK ) \arrow[r,symbol=\supseteq] & \Gamma_f \arrow[r] \arrow[d] & \PP^1(\OK ) \\
            & \X \arrow[ru, "f"', dotted]                  &
        \end{tikzcd}
    \end{center}
    The first projection $\X \times_{\OK } \PP^1(\OK ) \to \X$ induces a birational morphism $\Gamma_f \to \X$.
    The second projection induces a morphism $\Gamma_f \to \PP^1(\OK )$ which extends $f \colon \Gamma_f \dashrightarrow \PP^1(\OK )$.
    Let $\X' \to \Gamma_f$ denote the normalization of $\Gamma_f$.
    This gives a model of $\X'$ for which $f \colon \X' \to \PP^1(\OK )$ is a morphism.
    By Lemma~\ref{lem:condition_induced}~\ref{lem:condition_induced_ii}, this implies that $v \in V(\X')$.
\end{proof}


\subsection{Proof of the Second Part of Theorem~\ref{thm:bijectionModelsAndComponents}}\label{subsec:proof-of-the-second-part}
We fix a model $\X$ of $X$ and a nonempty subset $V \subset V(\X)$.
Denote by $\EE$ the set of irreducible (integral) components of $\X_s$ that induce the valuations in $V(\X)\setminus V$.

\begin{definition}
    A model $\X'$ together with a modification $\varphi \colon \X \to \X'$ such that for every integral vertical curve $E$ on $\X$,
    the set $\varphi(E)$ is a point if and only if $E \in \EE$ is called a \textit{contraction} of the $E \in \EE$.
\end{definition}

It can be shown that if a contraction exists, it is unique up to unique isomorphism.
Our goal now is to show that a contraction of $\EE$ exists.
Once this is done, the proof is complete, since for the contraction $\varphi \colon \X \to \X'$, we have by construction $V(\X') = V$.

\begin{lemma}[{{\cite[Lemma 8.3.29, Lemma 8.3.3]{Liu}}}] \label{lem:criterion_point}
    Let $\LL$ be an invertible sheaf on $\X$ generated by global sections $s_0, \dots, s_n$.
    Consider the morphism $f \colon \X \to \PP^n(\OK )$ associated to these sections.
    Let $E$ be a vertical curve of $\X_s$ with the reduced induced closed subscheme structure.
    Then $f(E)$ is reduced to a point if and only if $\LL|_E \cong \OO_E$.
\end{lemma}

\begin{proposition}[{{\cite[Proposition 8.3.30]{Liu}}}] \label{prop:exist_contr}
    The following are equivalent.
    \begin{enumerate}[(a)]
    \item The contraction $\varphi \colon \X \to \X'$ of $\EE$ exists. \label{prop:exist_contr_i}
    \item There exists a Cartier divisor $D$ on $\X$ such that $\deg(D|_{\X_g}) > 0$, that $\LL \cong \OO_\X(D)$ is generated
    by its global sections, and that for any integral vertical curve $E$, we have $\LL|_E \simeq \OO_E$ if and only if $E \in \EE$. \label{prop:exist_contr_ii}
    \end{enumerate}
\end{proposition}

\begin{proof}
    We only sketch (\ref{prop:exist_contr_ii}) $\implies$ (\ref{prop:exist_contr_i}).
    For the other implication and details see~\cite[Proposition 8.3.30]{Liu}.

    As $\X_g$ is an integral projective curve over a field, $\LL|_{\X_g}$
    is an ample divisor (\cite[Proposition 7.5.5]{Liu}). Replacing $D$ by a positive multiple (which
    does not change the set $E$), we can suppose that $\LL|_{\X_g}$ is very ample.

    Let $f \colon \X \to \PP^n(\OK )$ be a morphism associated to sections $s_0, \ldots, s_n \in \Gamma(\X, \LL)$ which generate $\LL$.
    Let $Z$ be the closed subset $f(\X) \subset \PP^n(\OK )$
    endowed with the reduced (hence integral) scheme structure.
    Then $f$ induces a dominant morphism $\X \to Z$ that we still denote by $f$.
    This is an isomorphism on the generic fibre because $\LL|_{\X_g}$ is very ample.
    In particular, $f \colon X \to Z$ is birational.
    Denote by $\pi \colon \X' \to Z$ the normalization morphism.
    As $\X$ is normal, $f$ uniquely factors into $\varphi \colon \X \to \X'$ followed by $\pi \colon \X' \to Z$.
    \begin{center}
        % https://tikzcd.yichuanshen.de/#N4Igdg9gJgpgziAXAbVABwnAlgFyxMJZARgBpiBdUkANwEMAbAVxiRAC0QBfU9TXfIRRkADFVqMWbADrSAGgHJuvEBmx4CREeXH1mrRCFlzu4mFADm8IqABmAJwgBbJACZqOCEm0T9bWyDUDHQARjAMAAr8GkIg9lgWABY4ynaOLojuIJ5IZL5ShrL09miJWKkgDs65Hl6IPnoFRtJo5VwUXEA
        \begin{tikzcd}
            & \X' \arrow[d, "\pi"] \\
            \X \arrow[r, "f"'] \arrow[ru, "\varphi"] & Z
        \end{tikzcd}
    \end{center}
    Let $E$ be an integral vertical curve on $\X$. Then we have
    \[
        f(E) \textrm{ is a point} \iff \LL|_E \simeq \mathcal{O}_E \iff E \in \EE
    \]
    where the first equivalence is a consequence of Lemma~\ref{lem:criterion_point} and the second one is true by assumption.
    We claim that $\pi \colon \X' \to Z$ is a finite morphism.
    This then completes the proof because $f(E)$ is reduced to a point if and only if $\varphi(E)$ is a point.
    The morphism $f$ is projective (\cite[Corollary 3.3.32(e)]{Liu}), so $\AA \cong f_*\OO_\X$ is a coherent
    sheaf on $Z$ (\cite[Corollary 5.3.5]{Liu}). Furthermore, $\AA$ is a sheaf of integrally closed algebras because $\X$ is normal.
    As a result, the normalization morphism $\pi \colon \X' \to Z$
    coincides with the canonical morphism $\SpecRel \AA \to Z$ and is a finite morphism (\cite[Exercise 5.1.17]{Liu}).
\end{proof}

For the construction of a suitable Cartier divisor, we furthermore need the following technical lemma.

\begin{lemma}[{{\cite[Lemma 8.3.32]{Liu}}}] \label{lem:global_sec}
    Let $D$ be an effective horizontal Cartier divisor on $\X$, and $\LL = \OO_\X(D)$.
    Then there exists a $m_0 \geq 1$ such that $\LL^{\otimes m}$ is generated by its global sections for every $m \geq m_0$.
\end{lemma}

Proposition~\ref{prop:exist_contr} and Lemma~\ref{lem:global_sec} show that for a contraction morphism to exist,
it suffices that there exists an effective Cartier divisor that meets the closed fibres in a suitable way.
To construct such a divisor, it is important that the ring $\OK $ is complete, which implies it is Henselian.

\begin{definition} %[{{\cite[8.3.33]{Liu}}}, {{\cite[\href{https://stacks.math.columbia.edu/tag/04GG}{Lemma 04GG}]{Stacks}}}]
    Let $A$ be a Noetherian local ring with maximal ideal $\m$ and residue field $k \coloneqq A/\m$.
    We say that $A$ is Henselian if one of the following equivalent conditions is fulfilled:
    \begin{itemize}
        \item Hensel's Lemma holds: For every monic $f \in A[T]$, and every simple root $a_0 \in k$ of $\overline{f} \in k[T]$,
            there exists a lift $a \in A$ (so $\overline{a} = a_0$) such that $f(a) = 0$.
        \item Every finite $A$-algebra is a direct sum of local $A$-algebras.
    \end{itemize}
\end{definition}

\begin{lemma}[{{\cite[8.3.35]{Liu}}}] \label{lem:exist_div}
    Let $x \in \X_s$ be a closed point on the special fiber, then the following assertions hold.
    \begin{enumerate}[(a)]
        \item There exists an irreducible horizontal Weil divisor $\Delta$ containing $x$. \label{lem:exist_div_i}
        \item There exists an effective horizontal Cartier divisor $D$ such that $\X_s \cap \Supp{D} = \{x\}$. \label{lem:exist_div_ii}
    \end{enumerate}
\end{lemma}

\begin{proof}
    \begin{enumerate}[(a)]
        \item Let $\m_x$ be the maximal ideal of $\OO_{\X,x}$ and denote by $\pi$ an uniformizer of $\OO_K$.
            Let $ \p_1, \dots, \p_n $ be the prime ideals of $\OO_{\X,x}$ that are minimal among those containing $\pi \OO_{\X,x}$.
            As $\dim \OO_{\X,x} = 2$ (\cite[Proposition 8.3.4]{Liu}), Krull's principal ideal theorem implies that $\m_x \not= \p_i$.
            There therefore exists an
            \[
                f \in \m_x \setminus (\cup_i \p_i).
            \]
            Thus $f$ is a regular element and $V(f)$ does not contain any irreducible component of $\Spec \OO_{\X_s, x}$.

            There exists an affine open subset $W \ni x$ and a regular element $g \in \OO_\X(W)$ such that $g_x = f$,
            and that $V(g) \cap W_s = \{x\}$.
            Let $\Delta$ be the closure in $\X$ of an irreducible component of $V(g)$ passing through $x$.
            Then $\Delta$ is an irreducible horizontal divisor passing through $x$. \label{lem:proof:exist_div_i}
        \item Let $E \coloneqq V(g)$ and let $W$ be as in~\ref{lem:proof:exist_div_i}.
            Then $E$ can be seen as an effective Cartier divisor on $W$.
            Moreover, $E \to \Spec \OO_K$ is quasi-projective and quasi-finite because $\dim E_s < \dim \X_s = 1$.
            By a corollary of Zariski's Main theorem (\cite[Corollary 4.4.8]{Liu}), there exists an affine open neighborhood
            $U$ of $x$ such that $E \cap U$ is an open subscheme of a scheme $Z$ that is finite over $\Spec \OO_K$.

            \begin{center}
                % https://tikzcd.yichuanshen.de/#N4Igdg9gJgpgziAXAbVABwnAlgFyxMJZABgBpiBdUkANwEMAbAVxiRAFEACAHW4GM6aTgFUQAX1LpMufIRQBGclVqMWbAFrjJIDNjwEiAJiXV6zVohC8Aymhh8e3APJOA+gGlxymFADm8IlAAMwAnCABbJDIQHAgkRRVzNl4cGAAPHBDw4Ag7ME4YcIAjHygsMF8xEGocOiwGNgALCAgAay1gsMjEBNikY0S1SxT0nGAg8twYKrEKMSA
                % https://tikzcd.yichuanshen.de/#N4Igdg9gJgpgziAXAbVABwnAlgFyxMJZAJgBoAGAXVJADcBDAGwFcYkQAtEAX1PU1z5CKAMwVqdJq3YAdGQGU0MAMYACOQHkNAfQDSPPiAzY8BIuXE0GLNohABRdTOX00qgKo8JMKAHN4RKAAZgBOEAC2SBYgOBBIAIxWUrYgcjgwAB44wEFYYLgw3AbBYZGIZDFxiNHW0nZpmTgh4cAQSmCqMOEARj5Qeb5FNDj0WIzsABYQEADWXtxAA
                \begin{tikzcd}
                    E \cap U \arrow[rr, "\textrm{open emb.}", hook] &  & Z \arrow[r, "\text{finite}"] & \Spec \OO_K
                \end{tikzcd}
            \end{center}

            Let $D$ be the connected component of $Z$ containing $x$.
            Since $\OO_{K}$ is a Henselian local ring, $Z$ is a disjoint union of local schemes.
            Hence $D$ is local, and consequently $D \subset E \cap U \subset \X$. %(Tippfehler in Liu)

            Finally, since $D$ is closed in $Z$, it is finite over $\Spec \OO_K$.
            A finite morphism $D \to \Spec \OO_K$ between affine schemes is projective (\cite[Exercise 3.3.22]{Liu}),
            and therefore proper (\cite[Theorem 3.3.30]{Liu}).
            So the morphisms
            \[
                D \to \X \to \Spec \OO_K \quad
                \textrm{and} \quad
                \X \to \OO_K
            \]
            are proper, which implies that $D \to \X$ is also proper (\cite[Proposition 3.3.16(e)]{Liu}). Consequently, $D$ is closed in $\X$ and therefore a Cartier divisor, as required.
    \end{enumerate}
\end{proof}

Now we can finally prove Proposition~\ref{prop:constructing_models}~\ref{prop:constructing_models_ii}.

\begin{proof}[Proof of Proposition \ref{prop:constructing_models}~(\ref{prop:constructing_models_ii})]
    Let $Z_1, \dots, Z_n$ be the irreducible components of $\X_s$ that induce the valuations in $V \subset V(\X)$.
    By Lemma~\ref{lem:exist_div}~\ref{lem:exist_div_ii}, there exists an effective horizontal Cartier divisor $D_i$
    such that $\X_s \cap \Supp D_i$ is a point of $Z_i$ that belongs to none of the components of $\EE$.
    Let $D = \sum_{i=1}^n D_i$.
    By virtue of Lemma~\ref{lem:global_sec}, if necessary replacing $D$ by a multiple, we can even suppose that
    $\LL \coloneqq \OO_\X(D)$ is generated by its global sections.

    The theorem then results from Proposition~\ref{prop:exist_contr}.
    Indeed, \(\LL|_E \simeq \OO_E\) if $E \in \EE$ since $E \cap \Supp D = \varnothing$;
    if $E \notin \EE$, then $E$ is some $Z_i$ and $\deg \LL|_E \geq \deg \OO_\X(D_i)|_{Z_i} > 0$, and hence \(\LL|_E \not\simeq \OO_E\).
\end{proof}

\begin{figure}[h!] \label{fig:illustration-of-proof-of-prop-constructing-models}
    \centering
    \tikzfig{figures/fig1}
    \caption{Illustration of the proof of Proposition \ref{prop:constructing_models}~(\ref{prop:constructing_models_ii})}
\end{figure}

\begin{remark}
    The proof of Proposition~\ref{prop:constructing_models}~\ref{prop:constructing_models_ii} shows that actually all models
    $\X$ of $X$ are projective.
    Indeed one just has to choose $V = V(\X)$ in the proof to obtain an effective Cartier divisor $D$ which meets
    all the irreducible components of the fiber $\X_s$.
    Since $\X_s$ is a projective curve (\cite[Exercise 7.5.4]{Liu}), $\OO_\X(D)|_{\X_s}$ is ample (\cite[Proposition 7.5.5]{Liu}).
    It follows from \cite[Corollary 5.3.24]{Liu} that $\OO_\X(D)$ is ample.
    So $\X \to \Spec \OO_K$ is projective (\cite[Proposition 5.3.6]{Liu}).
\end{remark}


\section{Permanent Models}\label{sec:permament-models}

We want to understand the behaviour of models under base change with respect to an extension of the base field.
Usually, $L|K$ will denote a finite field extension.
Since $v_K$ is Henselian, it has a unique extension to $v_L$ to $L$, and then $(L,v_L)$ also satisfies Assumption \ref{ass:v_K}.
Given a model $\X$ of $X$, we ask ourselves whether the base change $\X\times_{\OO_K}\OO_L$ is again a model of $X_L \coloneqq X\times_K L$.
Well, almost: all properties required of a model from Definition \ref{def:model} hold except possibly normality.
We obtain a model of $X_L$, called the {\em normalised base change} of $\X$ relative to $L|K$, by normalizing $\X\times_{\OO_K}\OO_L$.\footnote{
    Here we use the fact that Assumption \ref{ass:v_K} implies that $\OO_K$ is an \emph{excellent} ring.
    It follows that $\X\times_{\OO_K}\OO_L$ is an excellent scheme, and therefore its normalization is a finite morphism.
    See \cite[Chapter 8, in particular Theorem 8.2.39 (d)]{Liu}.
}
We say that a model $\X$ is \textit{permanent} if $\X \times_{\OK}\OL$ is itself a model of $X$:

\begin{definition} \label{def:permanent-model}
    A model $\X$ of $X$ is called \textit{permanent}, if $\X \times_{\OK} \OL$ is normal for all finite field extensions $L|K$.
\end{definition}

We have the following theorem about permanence:

\begin{theorem} \label{thm:permanence}
    \begin{enumerate}[(i)]
        \item For a model $\X$ of $X$,  the following three statements are equivalent:
            \begin{enumerate}[(a)]
                \item $\X$ is permanent,
                \item the special fiber $\X_s$ of $\X$ is reduced,
                \item all $v\in V(\X)$ are \emph{weakly ramified} over $v_K$, i.e.\ $v$ and $v_K$ have the same value group.
            \end{enumerate}
        \item Given a model $\X$, there exists a finite extension $L/K$ such that the normalized base change of $\X$ relative to $L|K$ is permanent.
    \end{enumerate}
\end{theorem}

In Subsection \ref{subsec:models-with-reduced-special-fibre}, we will give a proof of Part (i).
In Subsection \ref{subsec:models-under-base-field-extensions}, we will prove Part (ii) in the special case where all mulitplicities of $\X_s$
are prime to the characteristic of the residue field $k$. \textcolor{red}{TBD.}
In the last two sections of the chapter, we see proofs of Part (ii) in other important special cases.

\subsection{Models with Reduced Special Fibre}\label{subsec:models-with-reduced-special-fibre}

As in the last two sections, we assume $X$ to be a smooth, projective and absolutely irreducible curve over the local
field $K$.


\begin{remark}
    Since flatness and properness are preserved under base change, a model $\X$ of $X$ is permanent if and only if
    $\X \times_{\OK} \OL$ is a model of $X_L = X \times_K L$.
\end{remark}

\begin{definition}
    A \textit{pre-model} of $X$ is a flat and proper $\OK$-scheme $\X$ with $\X \times_{\OK} K = X$.
    That is, $\X$ satisfies all the conditions of a model except of being normal.
\end{definition}

\begin{lemma}[{\cite[Lemma 1.3]{ArzdorfWewers}}]\label{lem:conditions-for-normality}
    Let $\X$ be a pre-model of $X$.
    Then $\X$ is normal (i.e.\ a model) if and only if the following two conditions hold.
    \begin{enumerate}[(i)]
        \item The special fibre $\X_s$ has no embedded points, i.e.\ no point $x\in \X_s$ \label{lem:conditions-for-normality-i}
            satisfying $\m_x = \Ann{R}{r}$ for some $r\in R \coloneqq \OO_{\X_s,x}$ (see~\cite[Definition 7.1.6]{Liu}).
        \item The local ring $\OO_{\X,\eta}$ is a discrete valuation ring for every generic point $\eta \in \X_s$. \label{lem:conditions-for-normality-ii}
    \end{enumerate}
\end{lemma}

\begin{proof}
    According to Serre's criterion for normality~\cite[Theorem 8.2.23]{Liu},
    we need to check the following properties:
    \begin{itemize}
        \item[$(\mathrm{R}_1)$] $\X$ is regular at $x$ when $\codim{x} \leq 1$, and
        \item[$(\mathrm{S}_2)$] $\depth{\OO_{\X,x}} \geq 2$
    \end{itemize}
    for all $x\in \X$. \\
    The points on $\X$ with codimension $1$ are precisely the points on the generic fibre, which is regular by assumption,
    and the generic point on the special fibre.
    Since any discrete valuation ring is a local Noetherian domain and vice versa, $(\mathrm{R}_1)$ is
    equivalent to (\ref{lem:conditions-for-normality-ii}).
    By~\cite[Proposition 8.2.11]{Liu}, condition $(\mathrm{S}_2)$ is equivalent to say
    that $\depth{\OO_{\X_s, x}} \geq 1$ for all $x \in \X_s$.
    But this is always true for the generic point, and if $x$ is a different point it is equivalent to say that $x$ is not an embedded point.
\end{proof}

\begin{example}
    \begin{enumerate}
        \item Consider a pre-model with special fibre
            \begin{align}
                \X_s \coloneqq \Spec{k[x,y]/(x^2,xy)},
            \end{align}
            which is topologically equivalent to $\AA^1(k)$.
            The point corresponding to the maximal ideal $(x,y)$ is an embedded point.
        \item \textcolor{red}{\cite[Example 9.8.4]{Hartshorne}}.
    \end{enumerate}
\end{example}

\begin{proposition}[{\cite[Proposition 2.3]{ArzdorfWewers}}]
    If $\X_s$ is reduced, then $\X$ is permanent.
\end{proposition}

\begin{proof}
    Let $L|K$ be a finite extension and $l|k$ the extension of the corresponding residue fields.
    The special fibre of $\X' \coloneqq \X_{\OL} = \X \times_{\OK} \OL$ is $\X'_s := \X_s \times_k l$,
    which is reduced since $k$ is perfect by assumption.
    This means $\X'_s$ satisfies condition (ii) from Lemma~\ref{lem:conditions-for-normality}. \\
    For condition (i) let $\eta \in \X'_s$ be any generic point.
    The chalk $\OO_{\X', \eta}$ then is a Noetherian local domain of dimension $1$ and it holds
    \begin{align}
        \OO_{\X'_s, \eta} = \OO_{\X', \eta}/\pi_L\OO_{\X', \eta},
    \end{align}
    where $\pi_L \in \OL$ is a uniformising parameter.
    Since $\X'_s$ is reduced, $\pi_L\OO_{\X',\eta}$ is maximal ideal which means that
    condition (i) from Lemma~\ref{lem:conditions-for-normality} is satisfied as well.
    Thus $\X'$ is normal and hence $\X$ is permanent.
\end{proof}

\begin{proposition}[{\cite[Corollary 2.2 (ii)]{ArzdorfWewers}}]\label{prop:criterion-reducedness}
    Let $\X$ be a model of $X$, then the following statements are equivalent.
    \begin{enumerate}[(i)]
        \item $\X_s$ is reduced.
        \item $\OO_{\X_s, \eta}$ is a field for every generic point $\eta \in \X_s$.
        \item Every $v_Z \in V(\X)$ is weakly ramified.
    \end{enumerate}
\end{proposition}

\begin{proof}
    Since $\X$ is normal, there are no embedded points in $\X_s$ by Lemma~\ref{lem:conditions-for-normality}.
    Since $\X_s$ has dimension one, it is reduced if and only if $\X_s$ is reduced in codimension zero
    which in turn is the case if $\OO_{\X_s, \eta}$ is a field for every generic point $\eta \in \X$. \\
    For the second equivalence recall that according to Proposition~\ref{prop:multiplicities-coincide},
    for any generic point $\eta\in \X_s$ and $Z = \overline{\lbrace\eta\rbrace}$ it holds
    \begin{align}
        [\Gamma_{v_Z}: \Gamma_K] = \length \OO_{\X_s, \eta}.
    \end{align}
    This implies that $\OO_{\X_s, \eta}$ is a field if and only $v_Z$ is weakly ramified for all $v_Z \in V(\X)$.
\end{proof}

\begin{example}[Curve with reduced special fibre]
    Let $K = \QQ_3$ and consider the curve
    \begin{align}
        X: \, y^3 = 3x^2 + 1.
    \end{align}
    The equation is already defined over $\OK = \ZZ_3$, i.e.\ we can define a model with the same equation.
    The reduction $\X_s/\FF_3$ is given by $0 = \bar{y}^3 - 1 = (\bar{y} - 1)^3$ which means $\X_s$ is not reduced. \\
    The irreducible component is given by $Z: \, y-1$; denote the corresponding valuation by $v$.
    For $z \coloneqq y - 1$ we have
    \begin{align*}
        3x^2 &= y^3 - 1 \\
            &= (z+1)^3 - 1 \\
            &= z^3 + 3z^2 + 3z,
    \end{align*}
    where the right hand side is an Eisenstein polynomial.
    This results $v(z) = \frac{1}{3}v_G(3x^2) = \frac{1}{3} \notin \ZZ = \Gamma_K$.
    Let us extend the ground field to $L = \QQ_3[\theta]$ where $\theta^3 - 3 = 0$ and apply the coordinate transformation
    $(x,y) \mapsto (x_1, 1+ \theta y_1)$.
    We then obtain a birational curve
    \begin{align}
        X_1: \, y_1^3 + \theta^2 y_1^2 + \theta y_1 = x_1^2.
    \end{align}
    Since $y_1^3 \equiv x_1^2$ (mod $\theta$), $\X_s$ is reduced.
    It holds $v_G(\theta y_1) = v_G(y - 1) = \frac{1}{3} = v_L(\theta) \in \Gamma_L$, especially $v_G(y_1) = 0$.
\end{example}


\subsection{Models under Base Field Extensions}\label{subsec:models-under-base-field-extensions}

Let $L|K$ be a finite field extension and let $\X$ be a model of $X$.
Denote by $\X' \coloneqq (X\times_{\OK} \OL)^\sim$ the normalised base change of $\X$ to $\OL$.
Let $\eta \in \X_s$ be a generic point and let $\eta' \in \X_s'$ be a generic point lying over $\eta$.
Since $\X$ (resp. $\X'$) is normal we have that $R \coloneqq \OO_{\X, \eta}$ (resp. $ \coloneqq \OO_{\X', \eta'}$)
is a discrete valuation ring dominating $\OK$ (resp. $\OL$) by
Lemma~\ref{lem:conditions-for-normality} (\ref{lem:conditions-for-normality-i}), that is
\begin{itemize}
    \item $R \supset \OO_K$ (resp. $S \supset \OL$), and
    \item $\m_R \cap \OK = \pi_K\OK$ (resp. $\m_S \cap \OL = \pi_L\OL$).
\end{itemize}
We arrive at the following diagram of fields.


\begin{figure}[H]\label{fig:valuations-and-field-extensions}
    \centering
    \resizebox{0.5\textwidth}{!}{
        \begin{tikzcd}[ampersand replacement=\&]
            {}
            \& {(F_{X_L},\{w_{i,1},\dots,w_{i,j_i}\})} \\
            {(F_X,\{v_i\})} \arrow[ru, no head]
            \& {} \\
            {}
            \& {(L(x),v_{{\scriptscriptstyle G}, L})} \arrow[uu, no head]  \\
            {(K(x),v_{{\scriptscriptstyle G}, K})} \arrow[ru, no head] \arrow[uu, no head]
            \& {} \\
            {}
            \& {(L,v_L)} \arrow[uu, no head]  \\
            {(K,v_K)} \arrow[ru, no head] \arrow[uu, no head]
            \& {} \\
        \end{tikzcd}
    }
\end{figure}

Here the $v_i$ are the valuations on $F_X$ corresponding to the generic points of the components of $\X_s$.
The transcendental element of the field extension $K(x)|K$ can be chosen such that the Gauß valuation $v_{{\scriptscriptstyle G}, K}$
on $K(x)$ is compatible with the extension of valuations. \\
$w_{i, j}$ represent the valuations corresponding to the irreducible component of $\X_s'$ whose generic point lies over
the generic point of $Z_i \subset \X_s$ which corresponds to $v_i$.

\begin{theorem}[{\cite[Proposition 2.3]{ArzdorfWewers}}]
    Let $\X$ be a model of $X$.
    Then there is a finite extension $L|K$ such that $\X' = (\X \times_{\OK}\OL)^\sim$ has a reduced special fibre.
\end{theorem}

\begin{proof}
    By~\cite[Theorem 1.0]{Epp} there exists a field extension $L|K$ such that $w_{i,j}$ is weakly unramified over $v_L$.
    According to Proposition~\ref{prop:criterion-reducedness}, the special fibre of $\X'$ is reduced.
\end{proof}

\begin{lemma}[Abhyanka]\label{lem:abhyanka}
    Let $L|K$ be a finite extension.
    Let $\lbrace v_1, \dots, v_n \rbrace = V(X)$ denote the extended valuation on the function field $F_X|K$.
    We have that $F_{X_L} = F_X L$ and with the extended valuation denoted as $w_{1,1}, \dots, w_{n, j_n}$ respectively.
    Assume that $v_i|v_{{\scriptscriptstyle G}, K}$ are tamely ramified for all $i$.
    If $e_i \coloneqq e(v_i|v_{{\scriptscriptstyle G},K})$ divides
    $e(v_{{\scriptscriptstyle G}, L}, v_{{\scriptscriptstyle G}, K}) = e(v_L|v_K)$ for all $i$, then $w_{i,j} \in V(X_L)$
    is weakly unramified over $v_{{\scriptscriptstyle G}, L}$ for all $i,j$.
\end{lemma}

\begin{proof}
    We will use the equivalence of Proposition~\ref{prop:criterion-reducedness} without further mention,
    that is the equivalence of the special fibre being reduced and the valuations being mildly ramified.
    Without loss of generality we may assume the following assertions:
    \begin{itemize}
        \item $n = 1$ and $j_1 = 1$ since the argument can be constructed inductively,
        \item $e \coloneqq e_1 = e(v_1|v_{{\scriptscriptstyle G},K}) = e(v_{{\scriptscriptstyle G},L}|v_{{\scriptscriptstyle G},K})$
            since we assume the residue field to be perfect, and
        \item $K$ contains the $e$-th root of unity, i.e. $\pi_K = \pi_K^e$, since we can always replace $K$ by a suitable
            Kummer extension without violating the other assumptions and obtain a compatible diagram regarding the valuations.
    \end{itemize}
    Since $F_{X_L} = F_X L$ we have that $[\Gamma_{w_1}: \Gamma_{v_1}] = 1$.
    By multiplicativity of the ramification indices we conclude $[\Gamma_{w_1}: \Gamma_{v_L}] = 1$.
\end{proof}

\begin{example}[Reduction after wildly ramified field extension]
    Consider the curve
    \begin{align}
        X: \, y^2 = 8x^2 - 2
    \end{align}
    defined over the field $K \coloneqq \QQ_2$.
    This example illustrates that the right choice of the field extension is crucial for having a reduced special fibre,
    since a wild ramified field extension is not unique.
    \begin{enumerate}[(i)]
        \item First let us consider the extension $L_1 = K(\theta_1)$ where $\theta_1^2 - 2 = 0$ and the coordinate transformation
            $(x,y) \mapsto (x_1, 2y_1 - \theta_1 - 2)$.
            We obtain a birational curve
            \begin{align}
                X_1: \, y_1^2 - (\theta_1 + 4)y_1 = 2x_1^3 - (2 + \theta_1)
            \end{align}
            with $v_{{\scriptscriptstyle G}, L_1}(y_1) = \frac{1}{4}$.
            We have $[\Gamma_{F_{X_1}}: \Gamma_{L_1}] = 2 \neq 1$.
        \item Let us consider the field extension $L_2 = K(\theta_2)$ where $\theta_2^2 + 2 = 0$
            and the coordinate transformation $(x,y) \mapsto (x_2, 2\theta_2 y_2 + \theta_2)$.
            We obtain a birational curve
            \begin{align}
                X_2: \, y_2^2 + y_2 = -x_2^3
            \end{align}
            which is given by an equation of AS-form, i.e.\ its special fibre is a reduced curve.
            Thus $[\Gamma_{F_{X_2}}: \Gamma_{L_2}] = 1$.
    \end{enumerate}
\end{example}


\section{MacLane Valuations}

In his seminal work \cite{MacLane}, S.~MacLane considered and essentially solved the following problem.
Let $v_K$ be a discrete valuation on a field $K$.
What are the extensions of $v_K$ to an (arbitrary) pseudo-valuation $v$ on the polynomial ring $K[x]$?
How can we describe such pseudo-valuations explicitly?
These results have been applied to the problem of constructing models of curves in J.~Rüth's thesis, \cite{Rueth}. 
In this section, we give a brief overview.

\begin{definition} \label{def:pseudo-valuation}
    We call $v: K[x]: \to \hat{\RR}$ a pseudo-valuation if the following conditions are satisfied:
    \begin{itemize}
        \item $v(fg) = v(f) + v(g)$,
        \item $v(f + g) \geq \min\lbrace v(f), v(g) \rbrace$, and
        \item $v(0) = \infty$.
    \end{itemize}
\end{definition}

\begin{remark}
    Let $v$ be a pseudo-valuation on $K[x]$.
    \begin{enumerate} [(i)]
        \item $\p_v \coloneqq v^{-1}(\infty) < K[x]$ is a prime ideal; this is a direct consequence of the definition.
        \item Let $L|K$ be a finite field extension with minimal polynomial $\varphi \in K[x]$.
            Then we have a one-to-one correspondence between the valuations $v_L: L \to \hat{\RR}$ extending $v_K$ and 
            and the pseudo-valuations $v: K[x] \to \hat{\RR}$ with $v(\varphi)=\infty$.
            It is given by the mutually inverse maps
            \begin{align}
                v_L \mapsto \left( K[x] \twoheadrightarrow L \overset{v}{\to} \hat{\RR} \right)
            \end{align}
            and
            \begin{align}
                v \mapsto \left(K[x]/\p_v \overset{\bar{v}}{\to} \hat{\RR}\right).
            \end{align}
    \end{enumerate}
\end{remark}

Now let $V(K[x])$ be the set of all pseudo-valuations on $K[x]$ extending $v_K$.
Actually, this is the set underlying the $K$-analytic space $\an{\AA^1(K)}$ associated to the affine line over $K$.
